\providecommand{\kBKM}{\kappa_{\mathrm{BKM}}}
\providecommand{\kch}{\kappa_{\mathrm{ch}}}
\providecommand{\kcat}{\kappa_{\mathrm{cat}}}
\providecommand{\Sp}{\mathrm{Sp}}
\providecommand{\ZZ}{\mathbb{Z}}
\providecommand{\Mtf}{M_{24}}
\providecommand{\CM}{C_{\Mtf}}

\section{Five Proof Obligations for Conjecture~1$'$}
\label{sec:conj1-proof-obligations}

Conjecture~1$'$ concerns the automorphic, BKM, DT, and moonshine
realisations of the CHL forms attached to \(N\in\{1,2,3,4,6\}\) and to
the exceptional \(N\in\{5,7,8\}\) super-VOA range.  The assertions below
are the exact proof obligations left by the verified literature.

\subsection*{The Twisted Weight-Five Dimension}
\label{subsec:twisted-weight-five-dimension}

The existence of the Igusa square-root denominator comes from the
Gritsenko-Nikulin/Borcherds lift
\[
  \operatorname{Bor}(\phi_{0,1})=\Delta_5 .
\]
Uniqueness in the line
\[
  S_5(\Sp_4(\ZZ),\nu_{\Delta_5})
\]
requires a dimension computation for the sign-character space.  The
untwisted trace formula of Ibukiyama-Wakatsuki applies to
\(S_k(\Sp_4(\ZZ))\) with trivial character.  The Fourier-Jacobi
injection sends the twisted space into Jacobi forms with the
corresponding theta multiplier:
\[
  S_5(\Sp_4(\ZZ),\nu_{\Delta_5})
  \hookrightarrow
  \bigoplus_{m\ge 1}
  J^{\mathrm{cusp}}_{5,m}
  (\mathrm{SL}_2(\ZZ),\nu_{\Delta_5}^{1/2}) .
\]
The first coefficient lies in the half-integral theta-multiplier space,
so the comparison with the untwisted identity
\[
  J^{\mathrm{cusp}}_{5,1}(\mathrm{SL}_2(\ZZ))\cong
  S_6(\mathrm{SL}_2(\ZZ))=0
\]
proves no dimension statement for the sign-character space.  The
remaining obligation is a direct trace computation for
\(\nu_{\Delta_5}\), or a primary theorem proving
\(\dim S_5(\Sp_4(\ZZ),\nu_{\Delta_5})=1\).

\subsection*{Scalar Normalisation of the Four Constructions}
\label{subsec:scalar-normalisation}

If the twisted weight-five space is one-dimensional, the Arthur,
Kudla-Millson, Harvey-Moore, and Borcherds constructions land on a
common line:
\[
  F_i=C_i\Delta_5,\qquad i=1,2,3,4 .
\]
The equality \(C_1=C_2=C_3=C_4=1\) is four scalar normalisation
theorems.  The normalisations come from different data:
\begin{enumerate}[label=\textup{(\roman*)}]
\item Arthur packet multiplicities and endoscopic signs;
\item the Heegner divisor class and the Bruinier-Funke regulator in the
Kudla-Millson lift;
\item heterotic charge-lattice and GSO conventions in the
Harvey-Moore vacuum amplitude;
\item the Jacobi constant term in the Borcherds lift.
\end{enumerate}
Line uniqueness gives scalar proportionality.  Equality of the scalars
requires matching a single Fourier, divisor, or BPS coefficient in each
normalisation, using the same primitive denominator convention.

\subsection*{The DT Range Beyond the Primitive \(N=1\) Case}
\label{subsec:dt-range}

Oberdieck-Pixton prove the \(K3\times E\) reduced theory for primitive
K3-factor classes in the \(N=1\) setting.  Bryan-Oberdieck prove
selected primitive \(N=2\) base cases with square in \(\{-2,0\}\) by
topological-vertex methods.  The general assertions
\[
  Z^{\mathrm{DT}}_{[K3\times E/\ZZ_N]}
  =
  -\,\Phi_N^{-1},
  \qquad
  N\in\{2,3,4,6\},
\]
require new input: the reduced obstruction theory on the quotient, the
orientation line, the primitive-to-imprimitive divisor sums, and the
comparison with the CHL Borcherds product.  The verified literature
supplies the \(N=1\) primitive theorem and selected \(N=2\) checks; the
remaining CHL rows require their own DT proof.

\subsection*{Automorphic Product Versus BKM Denominator}
\label{subsec:automorphic-product-bkm}

Clery-Gritsenko construct paramodular Siegel cusp forms and write
explicit multiplicative products.  A BKM denominator theorem requires
additional root-theoretic data.  The product must determine a
Lorentzian lattice, a Weyl chamber, real simple roots, imaginary simple
roots, parity, and Borcherds-Serre relations with multiplicities
\[
  \mathrm{mult}(\alpha)
  =
  c_{g,h}\!\left(\frac{\alpha^2}{2},\alpha\cdot \ell\right).
\]
Borcherds' automorphic-product machinery supplies this conclusion under
its singular-weight and lattice hypotheses.  For the CHL weights
\[
  k_N=\frac{10}{N+1},\qquad N\in\{2,3,4,6\},
\]
those hypotheses must be checked for the actual lattice and character.
The automorphic form and its infinite product are established inputs;
the finite BKM recognition is a separate theorem.

\subsection*{The \(M_{24}\) Third-Cohomology Input}
\label{subsec:m24-third-cohomology}

The Mathieu cocycle input uses
\[
  H^3(M_{24},U(1))\cong\ZZ/12\oplus\ZZ/2 .
\]
Mason's 1994 paper computes the Schur multiplier and studies the quantum
double context.  Gaberdiel-Persson-Volpato cite the third cohomology
calculation to Thomas and to finite-group cohomology computations.  The
programme needs a reproducible computation of the displayed group and
of its restriction to the cyclic subgroups generated by the relevant
CHL elements.

For \(N=6\), restriction to \(C_6\) already shows the issue:
\[
  H^3(C_6,U(1))\cong\ZZ/6 .
\]
The \(\ZZ/12\)-summand has a nontrivial kernel under reduction modulo
six, and the additional \(\ZZ/2\)-summand is third cohomology.  The
perfectness identity \(H_1(M_{24},\ZZ)=0\) records degree-one
abelianisation and gives no third-cohomology vanishing.  Vanishing of
the K3 sigma-model cocycle on each CHL cyclic subgroup is therefore an
explicit restriction calculation.

\subsection*{Conclusion}
\label{subsec:conj1-proof-obligation-conclusion}

The \(N=1\) primitive automorphic denominator is supported by the
Gritsenko-Nikulin/Borcherds construction and the K3 reduced DT theorem.
The full Conjecture~1$'$ package requires five additional arguments:
the sign-character weight-five dimension, four scalar normalisations,
CHL quotient DT theorems for \(N\in\{2,3,4,6\}\), BKM recognition for
the paramodular products, and the \(M_{24}\) third-cohomology
restriction calculation.

\paragraph{Primary Sources.}
Ibukiyama-Wakatsuki, arXiv:1305.4886 (2014); Clery-Gritsenko,
\emph{Nagoya Math.\ J.}\ 199 (2010), arXiv:0812.3962; Gaberdiel,
Persson, and Volpato, arXiv:1211.7074 (2013); Mason,
\emph{Proc.\ London Math.\ Soc.}\ (1994); Gritsenko-Nikulin 1995,
1998; Borcherds, \emph{Invent.\ Math.}\ 120 (1995);
Oberdieck-Pixton, \emph{Invent.\ Math.}\ 222 (2020);
Bryan-Oberdieck, arXiv:1811.06102 (2018); Pitale-Schmidt,
\emph{JNT} 134 (2014); Taibi 2019.
